\section{Resultados}
\label{sec_resultados}

Esta seção apresenta a avaliação do serviço de metadados proposto. A avaliação é organizada em cinco partes: a metodologia experimental adotada, os testes funcionais que verificam a corretude do provisionamento, os testes de desempenho que quantificam a latência e a escalabilidade do serviço, a avaliação de alta disponibilidade baseada em consenso RAFT, e uma comparação de funcionalidades com serviços de metadados de provedores de nuvem públicos.

De modo geral, os experimentos confirmam a compatibilidade funcional do serviço com o \textit{cloud-init}: instâncias reais selecionam o \textit{datasource} de rede, obtêm seus metadados e concluem a inicialização com sucesso. As latências de resposta situam-se na faixa de dezenas de milissegundos e, ao contrário da hipótese inicial de degradação por serialização de escrita do SQLite, o serviço manteve latência estável e ausência de erros até 200 contêineres simultâneos. A avaliação de alta disponibilidade demonstrou formação de \textit{cluster}, replicação de estado e reeleição de líder da ordem de segundos após falha. Os detalhes de cada dimensão são apresentados a seguir.

\subsection{Metodologia de Avaliação}
\label{subsec_metodologia}

Os experimentos foram conduzidos em uma instância de nuvem Google Compute Engine do tipo \texttt{e2-standard-4}, com 4 vCPUs (processador Intel Xeon a 2{,}20~GHz), 15~GiB de memória RAM, executando Ubuntu 24.04 LTS (\textit{kernel} 6.17.0-1020-gcp) e Incus 6.0.0. O serviço de metadados foi implantado no mesmo host das instâncias.

Um aspecto prático relevante emergiu durante a configuração do ambiente: em provedores de nuvem pública (Google Cloud, AWS e Azure), o endereço \texttt{169.254.169.254} já é utilizado pelo próprio serviço de metadados do provedor, de modo que sua reutilização no host interfere na resolução de nomes e no acesso aos metadados da própria máquina virtual. Por essa razão, no ambiente de testes o \textit{cloud-init} foi apontado, via \textit{datasource} NoCloud, para o endereço de \textit{gateway} da \textit{bridge} do Incus (\texttt{10.10.10.1:8080}), diretamente alcançável pelas instâncias. Esse achado constitui, por si só, uma observação de relevância prática: a convenção do endereço reservado, embora adequada para instâncias hospedadas dentro de um provedor, não é diretamente transponível para um host que já opera sob tal provedor.

Para avaliar o comportamento sob diferentes cargas, foram criados grupos de instâncias com cardinalidades crescentes, de 5 até 200 contêineres simultâneos. Cada instância foi configurada com o \textit{cloud-init} habilitado, apontando para o serviço de metadados proposto como fonte de configuração. Durante a inicialização, o \textit{cloud-init} realizou requisições aos \textit{endpoints} do serviço (\texttt{/meta-data}, \texttt{/user-data}, \texttt{/vendor-data} e \texttt{/network-config}), seguindo o fluxo de \textit{datasource} documentado em \cite{cloudinit2024}.

As medições foram realizadas em três dimensões: (a) \textbf{latência de resposta da API}, correspondente ao tempo decorrido entre o envio da requisição HTTP pela instância e o recebimento da resposta; (b) \textbf{acurácia de sincronização}, verificando se os metadados armazenados refletem o estado atual das instâncias conforme reportado pelo Incus; e (c) \textbf{corretude da configuração}, avaliando se as configurações entregues foram efetivamente aplicadas pelas instâncias ao término da inicialização.

A coleta de latência foi realizada com a ferramenta \texttt{hey}, disparando 5000 requisições por configuração com concorrência de 50 requisições simultâneas a partir de um contêiner gerador de carga. Os resultados de latência são reportados nos percentis p50, p95 e p99, seguindo a prática adotada em sistemas de produção para caracterizar tanto o comportamento típico quanto os casos de cauda da distribuição.

\subsection{Avaliação Funcional}
\label{subsec_avaliacao_funcional}

A avaliação funcional verificou se o serviço de metadados entrega as informações corretas a cada instância e se essas informações são interpretadas e aplicadas com êxito pelo \textit{cloud-init}. As propriedades testadas e seus resultados são apresentados na Tabela~\ref{tab:funcional}.

\begin{enumerate}
    \item \textbf{Isolamento de metadados}: cada instância deve receber exclusivamente os metadados associados ao seu próprio identificador. Múltiplas instâncias realizaram requisições simultâneas ao serviço, e confirmou-se que o campo \texttt{local-hostname} correspondia unicamente à instância solicitante.

    \item \textbf{Isolamento sob falsificação de origem}: como o serviço identifica a instância pelo endereço IP de origem da requisição, foi testada a resistência à falsificação por meio do cabeçalho \texttt{X-Forwarded-For}. Confirmou-se que uma requisição forjada com o IP de outra instância não retornou os dados da vítima, validando a desativação da confiança em \textit{proxies} (\texttt{SetTrustedProxies(nil)}).

    \item \textbf{Execução de \textit{user-data}}: um \textit{script} de \textit{user-data} injetado via chave \texttt{cloud-init.user-data} do Incus foi executado durante a inicialização, criando um arquivo sentinela verificado após o término do \textit{cloud-init}.

    \item \textbf{Aplicação de configuração de rede}: a configuração entregue pelo \textit{endpoint} \texttt{/network-config} foi servida como YAML válido e capturada para inspeção do estado das interfaces via \texttt{ip addr}.

    \item \textbf{Atualização após mudança de estado}: após a reinicialização de uma instância, os metadados persistidos refletiram o novo estado dentro da janela de sincronização periódica.
\end{enumerate}

\begin{table}[ht]
\centering
\caption{Resultados dos testes funcionais.}
\label{tab:funcional}
\begin{tabular}{lc}
\hline
\textbf{Propriedade} & \textbf{Resultado} \\
\hline
Isolamento de metadados & Aprovado \\
Isolamento sob falsificação de IP (\texttt{X-Forwarded-For}) & Aprovado \\
Execução de \textit{user-data} & Aprovado \\
Entrega de \textit{network-config} & Aprovado \\
Propagação de estado após reinício & Aprovado \\
\hline
\end{tabular}
\end{table}

Todas as propriedades foram atendidas. Em particular, uma inicialização real de \textit{cloud-init} selecionou o \textit{datasource} \texttt{NoCloud} de rede, obteve seus metadados e \textit{user-data} do serviço (respostas HTTP 200) e concluiu com \texttt{cloud-init status: done}, evidenciando a compatibilidade funcional de ponta a ponta.

\subsection{Avaliação de Desempenho}
\label{subsec_avaliacao_desempenho}

A avaliação de desempenho concentrou-se em três métricas: latência de resposta da API, latência de sincronização e escalabilidade sob variação do número de instâncias.

\subsubsection{Latência de Resposta da API}

A latência foi medida para os quatro \textit{endpoints} do serviço. A Tabela~\ref{tab:latencia} apresenta os percentis p50, p95 e p99 obtidos com 5000 requisições por \textit{endpoint}. O \textit{endpoint} \texttt{/vendor-data} respondeu com código HTTP 404 por não haver dado de \textit{vendor} configurado no cenário --- comportamento correto, tolerado pelo \textit{cloud-init} ---, mas suas latências permanecem representativas do custo de processamento da requisição.

\begin{table}[ht]
\centering
\caption{Latência de resposta por \textit{endpoint} (5000 requisições, concorrência 50).}
\label{tab:latencia}
\begin{tabular}{lccc}
\hline
\textbf{Endpoint} & \textbf{p50 (ms)} & \textbf{p95 (ms)} & \textbf{p99 (ms)} \\
\hline
\texttt{/meta-data}       & 13{,}2 & 53{,}7 & 83{,}2 \\
\texttt{/user-data}       & 18{,}9 & 53{,}6 & 72{,}4 \\
\texttt{/network-config}  & 11{,}4 & 47{,}2 & 72{,}4 \\
\texttt{/vendor-data}     & 17{,}0 & 48{,}4 & 68{,}9 \\
\hline
\end{tabular}
\end{table}

As latências medianas situam-se em torno de 11--19~ms e o percentil p99 abaixo de 90~ms para todos os \textit{endpoints}, valores compatíveis com as janelas de tempo disponíveis durante a inicialização de instâncias.

\subsubsection{Latência de Sincronização}

A latência de sincronização --- intervalo entre a criação de uma instância no Incus e a disponibilização de seu registro de metadados --- foi medida em aproximadamente \textbf{20{,}5~s}. Esse valor é consistente com a janela do ciclo de sincronização periódica adotado (\textit{polling} a cada 10~s) somada ao tempo de inicialização de rede da instância, correspondendo tipicamente a cerca de dois ciclos de sincronização.

\subsubsection{Escalabilidade}

Para avaliar o comportamento com o aumento do número de instâncias simultâneas, mediu-se a latência de resposta do \textit{endpoint} \texttt{/meta-data} para configurações de 5 a 200 contêineres. A hipótese inicial era de que, por utilizar SQLite como mecanismo de persistência, o serviço apresentasse degradação de desempenho sob alta concorrência de escrita, dado que o SQLite serializa operações de escrita \cite{openstack_metadata}. A Tabela~\ref{tab:escalabilidade} apresenta os resultados.

\begin{table}[ht]
\centering
\caption{Escalabilidade: latência de \texttt{/meta-data} e memória por número de instâncias.}
\label{tab:escalabilidade}
\begin{tabular}{lcccccc}
\hline
\textbf{Instâncias} & \textbf{p50 (ms)} & \textbf{p95 (ms)} & \textbf{p99 (ms)} & \textbf{Erros} & \textbf{Mem. (MB)} \\
\hline
5   & 16{,}2 & 60{,}3 & 92{,}4 & 0{,}0\% & 1152 \\
10  & 13{,}1 & 52{,}7 & 78{,}2 & 0{,}0\% & 1325 \\
25  & 12{,}4 & 51{,}8 & 77{,}2 & 0{,}0\% & 1373 \\
50  & 12{,}9 & 52{,}0 & 80{,}8 & 0{,}0\% & 1733 \\
100 & 14{,}4 & 58{,}6 & 92{,}1 & 0{,}0\% & 1763 \\
150 & 12{,}8 & 53{,}4 & 82{,}2 & 0{,}0\% & 1638 \\
200 & 12{,}8 & 54{,}2 & 87{,}2 & 0{,}0\% & 1711 \\
\hline
\end{tabular}
\end{table}

Os dados \textbf{não corroboram} a hipótese de degradação por serialização de escrita nessa faixa de escala: a latência permaneceu estável (p50 entre 12 e 16~ms; p99 entre 77 e 98~ms) e a taxa de erro foi nula em todas as configurações, até 200 instâncias sobre 4 vCPUs. Esse resultado decorre da natureza do fluxo predominante --- requisições de leitura (\texttt{GET}) de metadados ---, enquanto as escritas de sincronização ocorrem apenas a cada 10~s e são absorvidas pelo modo WAL (\textit{Write-Ahead Logging}) e pelo \textit{timeout} de ocupação configurados no SQLite. Conclui-se que a serialização de escrita não constitui gargalo no perfil de carga avaliado; uma eventual degradação exigiria contagens de instâncias ou taxas de escrita substancialmente superiores, o que se registra como limite medido e não como tendência de degradação linear.

\subsection{Avaliação de Alta Disponibilidade}
\label{subsec_avaliacao_ha}

Para avaliar o mecanismo de alta disponibilidade baseado em consenso RAFT, provisionou-se um \textit{cluster} de três nós (\texttt{e2-standard-2}) com endereços internos fixos, um dos quais designado como nó de inicialização (\textit{bootstrap}). Adotou-se a topologia de \textbf{réplicas sobre um Incus compartilhado}: os três serviços de metadados apontam para a mesma instância do Incus, de modo que o líder sincroniza o estado a partir de uma visão única de instâncias e o replica aos seguidores via \textit{log} do RAFT. A Tabela~\ref{tab:ha} resume os resultados.

\begin{table}[ht]
\centering
\caption{Resultados da avaliação de alta disponibilidade (\textit{cluster} de 3 nós).}
\label{tab:ha}
\begin{tabular}{ll}
\hline
\textbf{Propriedade} & \textbf{Resultado} \\
\hline
Formação do \textit{cluster} & 1 líder + 2 seguidores \\
Replicação de estado & instância criada no líder presente nos 3 nós \\
Reeleição após falha do líder & 2{,}40~s \\
Disponibilidade dos dados após \textit{failover} & preservada (sem remoção indevida) \\
Reintegração de nó & nó reiniciado reingressa como seguidor \\
\hline
\end{tabular}
\end{table}

O \textit{cluster} formou-se corretamente, com eleição de um único líder. Uma instância criada no Incus compartilhado teve seu registro replicado aos três nós. Ao encerrar-se abruptamente o serviço do líder (mantendo-se o Incus compartilhado ativo), um novo líder foi eleito em \textbf{2{,}40~s} --- valor coerente com os \textit{timeouts} configurados (\textit{heartbeat} de 1~s e eleição de 1~s) ---, e os dados replicados permaneceram disponíveis nos nós sobreviventes, sem remoção indevida. Um nó previamente encerrado foi capaz de reingressar ao \textit{cluster} e recuperar o estado a partir do \textit{log}.

Observou-se que a corretude da rotina de reconciliação (que remove registros de instâncias ausentes do Incus) depende da topologia adotada: sob a topologia de Incus compartilhado, todo nó que possa assumir a liderança observa o mesmo conjunto de instâncias, de modo que a reconciliação é correta e o \textit{failover} não provoca perda de dados. Uma topologia alternativa, com um Incus por nó (adequada a cenários de borda distribuída), exigiria reconciliação restrita às instâncias de cada nó de origem, de forma que um nó jamais removesse instâncias pertencentes a outro.

\subsection{Comparação com Provedores de Nuvem}
\label{subsec_comparacao}

Para contextualizar as capacidades do serviço proposto, a Tabela~\ref{tab:comparacao_campos} compara os campos de metadados disponibilizados com os oferecidos pelo AWS Instance Metadata Service \cite{aws_imds}, pelo serviço de metadados do Google Cloud Platform \cite{gcp_metadata} e pelo serviço de metadados do OpenStack \cite{openstack_metadata}.

\begin{table}[ht]
\centering
\caption{Comparação de campos de metadados (\ding{51} suportado; \ding{55} não suportado; P: parcial).}
\label{tab:comparacao_campos}
\begin{tabular}{lcccc}
\hline
\textbf{Campo} & \textbf{Proposto} & \textbf{AWS} & \textbf{GCP} & \textbf{OpenStack} \\
\hline
\texttt{instance-id}        & \ding{51} & \ding{51} & \ding{51} & \ding{51} \\
\texttt{local-hostname}     & \ding{51} & \ding{51} & \ding{51} & \ding{51} \\
\texttt{local-ipv4}         & \ding{51} & \ding{51} & \ding{51} & \ding{51} \\
Chaves SSH públicas         & \ding{51} & \ding{51} & \ding{51} & \ding{51} \\
\texttt{availability-zone}  & \ding{51} & \ding{51} & \ding{51} & P \\
Tipo de instância           & \ding{55} & \ding{51} & \ding{51} & \ding{51} \\
\texttt{user-data}          & \ding{51} & \ding{51} & \ding{51} & \ding{51} \\
\texttt{vendor-data}        & \ding{51} & \ding{55} & \ding{55} & \ding{51} \\
\texttt{network-config}     & \ding{51} & P & P & \ding{51} \\
\hline
\end{tabular}
\end{table}

O serviço proposto cobre os campos essenciais para a operação dos módulos mais comuns do \textit{cloud-init} --- \texttt{set\_hostname}, \texttt{users-groups}, \texttt{write-files}, \texttt{runcmd} e configuração de rede ---, oferecendo também entrega de \textit{vendor-data} por instância com \textit{fallback} global. Campos específicos de provedores públicos, como tipo de instância e atributos de faturamento, não são expostos, por não se aplicarem ao contexto de \textit{homelab} e computação de borda visado.

\subsection{Discussão}
\label{subsec_discussao}

Os resultados obtidos permitem avaliar em que medida o serviço de metadados proposto atende aos requisitos de um ambiente Incus de pequena escala. Do ponto de vista funcional, o serviço demonstrou corretude no isolamento de metadados --- inclusive sob tentativa de falsificação de origem --- e na entrega de configurações efetivamente aplicadas pelo \textit{cloud-init}, validando a premissa de compatibilidade com o \textit{datasource} de rede \cite{cloudinit2024}. Do ponto de vista de desempenho, a natureza leve da implementação traduziu-se em latências de dezenas de milissegundos e em escalabilidade estável até 200 instâncias. Do ponto de vista de disponibilidade, o mecanismo de consenso demonstrou eleição e recuperação de líder da ordem de segundos.

Entre as limitações do serviço proposto, destacam-se:

\begin{itemize}
    \item \textbf{Escalabilidade do SQLite}: embora não se tenha observado degradação até 200 instâncias, o SQLite serializa operações de escrita, o que poderia tornar-se um gargalo sob taxas de criação e destruição de instâncias muito superiores às avaliadas. Essa característica é inerente à escolha por uma solução sem dependências externas e é aceitável no perfil de carga de \textit{homelab} e computação de borda de pequena escala.

    \item \textbf{Cobertura parcial de campos de metadados}: o serviço implementa o subconjunto de campos necessários aos módulos mais comuns do \textit{cloud-init}, mas não cobre todos os campos de provedores como AWS e GCP \cite{aws_imds, gcp_metadata}. Cargas que dependam de campos específicos podem requerer extensões.

    \item \textbf{Dependência da topologia na alta disponibilidade}: as garantias de disponibilidade sem perda de dados no \textit{failover} pressupõem a topologia de réplicas sobre um Incus compartilhado; cenários com um Incus por nó demandam reconciliação por nó de origem.
\end{itemize}

Quanto às ameaças à validade, os experimentos de desempenho e escalabilidade foram executados em host único, o que não captura efeitos de rede real entre \textit{hypervisor} e instâncias em topologias amplamente distribuídas; a avaliação de alta disponibilidade, contudo, já emprega um \textit{cluster} de três nós. As instâncias utilizadas são contêineres do tipo sistema (\textit{system containers}), que compartilham o \textit{kernel} do host, podendo apresentar latências de inicialização distintas das de máquinas virtuais completas. Por fim, a variabilidade do intervalo de sincronização periódica introduz uma janela de inconsistência cujos efeitos dependem da frequência de mudanças de estado no ambiente em produção. Em conjunto, os achados sustentam o objetivo definido na Seção~\ref{sec_introducao}: prover, para ambientes Incus, um serviço de metadados compatível com \textit{cloud-init}, leve e com suporte a alta disponibilidade.
